% lit-core-proofs.tex
% Fragment for the paper (no preamble). Cluster: prior work for certified
% coordinate grids, filtering, conditioned complexity and exact output.
% Contents:
%   (A) Proposition lc:cellwise -- the corrected-grid bound is exactly the
%       best cellwise edge-concave vertex bound; deterministic proof of the
%       interpolation lemma; interval-specific conditional bound.
%   (B) Proposition lc:necessity and Corollary lc:illconditioned -- width
%       alone cannot give a logarithmic accuracy rate unless P = NP, so a
%       conditioning parameter is necessary; NP-hard width-two families are
%       necessarily ill-conditioned.
%   (C) Remark lc:eth -- exponential dependence on the bag size is
%       unavoidable under ETH even when no growth premise is needed.
%   (D) A draft related-work subsection with citations to lit-core.bib.
% Cross-references to the main text: lem:height (rational height),
% thm:approx (certified approximation), eq:semiconcavity, eq:growth.
% Notation follows the shared paper notation: X = prod_i X_i, F, \OPT = F^*,
% L, g, kappa, grids G_i, widths w_i(v), corrections d_i(v), D, Q, beta,
% min-marginals m_i(v).

\subsection{The corrected grid as a cellwise edge-concave bound}\label{sec:lc-cellwise}

Fix a product subbox $B=\prod_i B_i$ of the continuous hull of $X$ in which
no coordinate is fixed, and for each $i$ a finite grid
$G_i\subseteq B_i\cap X_i$ that contains both endpoints of $B_i$. An
\emph{interval} of $G_i$ is a pair of consecutive nodes $[a,a^+]$. Its
\emph{effective width} is
\[
\Delta_i([a,a^+])=
\begin{cases}
0, & i\in I_Z \text{ and } a^+-a=1,\\
a^+-a, & \text{otherwise.}
\end{cases}
\]
For a node $v\in G_i$ let $w_i(v)$ be the largest effective width of an
interval of $G_i$ that has $v$ as an endpoint; since $|G_i|\ge2$ such an
interval exists. Put $d_i(v)=Lw_i(v)^2/8$, $D(y)=\sum_id_i(y_i)$,
$Q=F-D$, $\beta=\min\{Q(y):y\in\prod_iG_i\}$, and
$m_i(v)=\min\{Q(y):y\in\prod_kG_k,\ y_i=v\}$.

\begin{definition}[Cells and cell bounds]
A \emph{cell} is a product $C=\prod_iI_i$ in which every $I_i=[a_i,a_i^+]$
is an interval of $G_i$. Its vertex set is
$\operatorname{vert}(C)=\prod_i\{a_i,a_i^+\}$, and its \emph{cell bound} is
\[
\beta(C)=\min_{v\in\operatorname{vert}(C)}F(v)
-\frac L8\sum_{i}\Delta_i(I_i)^2 .
\]
\end{definition}

The quantity $\beta(C)$ is the vertex bound obtained from an
edge-concave underestimator: subtracting the separable quadratic
$\frac L2\sum_i(x_i-\mathrm{mid}\,I_i)^2$ makes $F$ concave along every
coordinate of $C$, so its minimum over $C$ is attained at a vertex
\cite{Tardella2004,MeyerFloudas2005,Hasan2018}. For a single box this is the
lower bound of Bajaj and Hasan \cite[Theorem~1]{BajajHasan2020}, and
$\frac L8\sum_i\Delta_i^2$ is the classical maximum separation of
$\alpha$BB underestimators with $\alpha=L/2$
\cite{MaranasFloudas1994,AdjimanEtAl1998}. The next proposition shows that
the node-separable correction $D$ computes the best of these bounds over
\emph{all} cells of the product grid at once.

\begin{proposition}[Cellwise form of the corrected grid bound]
\label{prop:lc-cellwise}
Assume \eqref{eq:semiconcavity}.
\begin{enumerate}
\item[(a)] For every cell $C$ and every $z\in C\cap X$, $F(z)\ge\beta(C)$.
Every $z\in B\cap X$ lies in some cell.
\item[(b)] $\beta=\min_C\beta(C)$, the minimum being over all cells. In
particular $\beta\le\min\{F(z):z\in B\cap X\}$.
\item[(c)] For a coordinate $i$ and an interval $I=[a,a^+]$ of $G_i$, put
$\beta_i(I)=\min\{\beta(C):C_i=I\}$. Then
\[
\beta_i(I)=\min_{v\in\{a,a^+\}}\bigl(m_i(v)+d_i(v)\bigr)
-\frac L8\Delta_i(I)^2
\;\ge\;\min\{m_i(a),m_i(a^+)\},
\]
and $F(z)\ge\beta_i(I)$ for every $z\in B\cap X$ with $z_i\in I$.
\end{enumerate}
\end{proposition}

\begin{proof}
(a) Let $C=\prod_i[a_i,a_i^+]$ and $z\in C\cap X$. Let
$J=\{i:\Delta_i(I_i)>0\}$, let $c_i=(a_i+a_i^+)/2$, and define
\[
\varphi(y)=F(y)-\frac L2\sum_{i\in J}(y_i-c_i)^2 .
\]
Then $\varphi\le F$. If $i\notin J$, then $i\in I_Z$ and $a_i^+=a_i+1$, so
$z_i\in\{a_i,a_i^+\}$ because $z_i$ is an integer in $[a_i,a_i+1]$. If
$i\in J$, then for fixed values of the other coordinates the function
\[
t\longmapsto\varphi(\ldots,t,\ldots)
=\Bigl(F(\ldots,t,\ldots)-\frac L2t^2\Bigr)+Lc_it+\text{const}
\]
is concave on $[a_i,a_i^+]$ by \eqref{eq:semiconcavity}, because the
interval lies in the continuous hull of $X_i$. A concave function on an
interval attains its minimum at an endpoint. Process the coordinates of
$J$ one at a time and replace $z_i$ by an endpoint of $[a_i,a_i^+]$ at
which this one-dimensional function is smallest. Each step keeps the point
in $C\cap X$, since grid nodes are feasible, and does not increase
$\varphi$. The result is a vertex $v$ of $C$ with $\varphi(v)\le\varphi(z)$.
For $i\in J$ we have $(v_i-c_i)^2=\Delta_i(I_i)^2/4$, and the terms with
$i\notin J$ have $\Delta_i(I_i)=0$. Hence
\[
F(z)\ge\varphi(z)\ge\varphi(v)=F(v)-\frac L8\sum_i\Delta_i(I_i)^2\ge\beta(C).
\]
Finally, $G_i$ contains both endpoints of $B_i$, so every $z_i\in B_i\cap X_i$
lies in some interval of $G_i$, and $z$ lies in the product of these
intervals.

(b) For a node $v\in G_i$ let $\mathcal I_i(v)$ be the nonempty set of
intervals of $G_i$ having $v$ as an endpoint. The pairs $(C,v)$ with
$v\in\operatorname{vert}(C)$ are exactly the pairs consisting of a grid
point $v\in\prod_iG_i$ and a choice $I_i\in\mathcal I_i(v_i)$ for every
$i$, and these choices are independent across coordinates. Therefore
\[
\min_C\beta(C)
=\min_{v\in\prod_iG_i}\Bigl[F(v)-\frac L8\sum_i
\max_{I_i\in\mathcal I_i(v_i)}\Delta_i(I_i)^2\Bigr]
=\min_{v}\bigl[F(v)-D(v)\bigr]=\beta .
\]
The last assertion follows from (a).

(c) The same separation, with the $i$-th interval fixed to $I$, gives
\[
\beta_i(I)=\min_{v\in\{a,a^+\}}\min\Bigl\{F(y)-\sum_{j\ne i}d_j(y_j):
y\in\textstyle\prod_kG_k,\ y_i=v\Bigr\}-\frac L8\Delta_i(I)^2 ,
\]
and the inner minimum equals $m_i(v)+d_i(v)$ because $Q(y)=F(y)-
\sum_{j\ne i}d_j(y_j)-d_i(v)$ when $y_i=v$. Since $I\in\mathcal I_i(v)$ for
$v\in\{a,a^+\}$, we have $d_i(v)\ge L\Delta_i(I)^2/8$, which gives the
inequality. If $z\in B\cap X$ and $z_i\in I$, choose the $i$-th interval of
the cell in (a) to be $I$; then $F(z)\ge\beta(C)\ge\beta_i(I)$.
\end{proof}

\begin{remark}[What the proposition adds]\label{rem:lc-cellwise}
Part (a) is a deterministic proof of the interpolation lemma that avoids
randomized rounding. Part (b) identifies the corrected grid minimum with
the best cellwise edge-concave bound. Computing $\min_C\beta(C)$ by
enumeration costs $\prod_i(|G_i|-1)$ cells with $2^n$ vertex evaluations
each, as in box-by-box branch and bound; because $D$ is a sum of unary
terms, the same number is the minimum of a factorized function on the
product grid and is obtained by one min-sum computation on the supplied
tree decomposition, with tables of size $\max_i|G_i|^p$. Part (c)
gives an interval-specific filtering test $\beta_i(I)\le U$ that is
never weaker than $\min\{m_i(a),m_i(a^+)\}\le U$ and costs nothing extra.
It keeps the two properties used later: an interval adjacent to a
coordinate of a minimizer $y$ of $Q$ survives (take the adjacent interval of
maximal effective width, for which $\beta_i(I)\le m_i(y_i)=\beta\le U$),
and every retained interval has an endpoint $v$ and a grid point $z$ with
$z_i=v$ and $Q(z)\le U$ (because $d_i(v)\ge L\Delta_i(I)^2/8$).
\end{remark}

\begin{remark}[Coordinate-specific curvature]\label{rem:lc-coordinate-L}
Proposition~\ref{prop:lc-cellwise} and its proof hold verbatim when
\eqref{eq:semiconcavity} is assumed with a separate constant $L_i\ge0$ for
each coordinate and $L$ is replaced by $L_i$ in $d_i$ and in the $i$-th term
of $\beta(C)$. For a quadratic one may take $L_i=\max\{H_{ii},0\}$. A
coordinate with $L_i=0$ then has zero correction, and the two-point grid
$G_i=\{l_i,u_i\}$ suffices for it: by coordinatewise concavity some optimal
solution has that coordinate at an endpoint, which is the endpoint
reduction used by Del Pia and Khajavirad \cite{DelPiaKhajavirad2026}, and
the rounding and gap estimates then involve only the other coordinates. The
growth analysis is unchanged if $L$ is replaced by $\max_iL_i$.
\end{remark}

\subsection{A conditioning parameter is necessary}\label{sec:lc-necessity}

The next proposition shows that the dependence of the running time on a
conditioning parameter cannot be removed by any algorithm, already for
continuous box QPs whose interaction graph has treewidth two.

\begin{proposition}[Width alone does not give a logarithmic accuracy rate]
\label{prop:lc-necessity}
Suppose an algorithm receives a rational continuous box QP
\[
\min\{x^{\mathsf T}Qx+c^{\mathsf T}x:\ x\in[0,1]^n\},
\]
with $Q$ symmetric, a tree decomposition of its interaction graph of width
at most two, and an integer $q\ge1$, and returns a rational number
$\tilde v$ with $|\tilde v-\OPT|\le2^{-q}$ in time polynomial in $I+q$.
Then $\mathrm P=\mathrm{NP}$.
\end{proposition}

\begin{proof}
By Theorem~3 of Del Pia and Khajavirad \cite{DelPiaKhajavirad2026}, the
following problem is NP-hard: given integral symmetric $Q$ and integral $c$
whose interaction graph has treewidth at most two, and a rational
$\gamma=r/s$ with $s\ge1$, decide whether $\OPT\le\gamma$. A tree
decomposition of width at most two can be computed in linear time
\cite{Bodlaender1996}. With $D=1$ and $C_Q=\max\{1,\max_{jk}|Q_{jk}|\}$,
Lemma~\ref{lem:height} (compare \cite{Vavasis1990}) shows that $\OPT$ is a
rational number whose reduced denominator is at most
$V=(2nC_Q)^{2n}$. Choose
\[
q=2n\lceil\log_2(2nC_Q)\rceil+\lceil\log_2s\rceil+2,
\]
so that $\varepsilon=2^{-q}\le1/(4Vs)$; $q$ is polynomial in the input
length. Run the algorithm and answer ``yes'' exactly when
$\tilde v\le\gamma+\varepsilon$.

If $\OPT\le\gamma$, then $\tilde v\le\OPT+\varepsilon\le\gamma+\varepsilon$.
If $\OPT>\gamma$, the two distinct rationals $\OPT$ and $\gamma$ have
denominators at most $V$ and $s$, so $\OPT-\gamma\ge1/(Vs)\ge4\varepsilon$
and $\tilde v\ge\OPT-\varepsilon\ge\gamma+3\varepsilon>\gamma+\varepsilon$.
Thus the answer is correct, and the total time is polynomial in the input
length. An NP-hard problem would then be solvable in polynomial time.
\end{proof}

The certified algorithm of Theorem~\ref{thm:approx} returns $U$ and $\beta$
with $\beta\le\OPT\le U\le\beta+2^{-q}$, so $\tilde v=U$ is admissible.
Inspection of its proof shows that, for fixed $p$, the factor $f(p,\kappa)$
is bounded by a polynomial in $\kappa$: the state cap is
$O(\sqrt\kappa\log(n+2))$ per coordinate, bag tables have at most the
$p$-th power of the cap as size, and all bit lengths are polynomial in $I+q$
and the cap.

\begin{corollary}[NP-hard width-two families are ill-conditioned]
\label{cor:lc-illconditioned}
Let $\pi$ be a polynomial. The problem of deciding $\OPT\le\gamma$ for
rational mixed-integer box QPs given with a tree decomposition of width at
most two is solvable in polynomial time under the promise that the
optimizer is unique and $\kappa\le\pi(I)$. Consequently, unless
$\mathrm P=\mathrm{NP}$, for every polynomial $\pi$ every polynomial-time
reduction proving NP-hardness of this problem (such as the reduction in
\cite[Theorem~3]{DelPiaKhajavirad2026}) produces some instances that have
several optimizers or satisfy $\kappa>\pi(I)$.
\end{corollary}

\begin{proof}
Run the algorithm of Theorem~\ref{thm:approx} with the accuracy $q$ chosen
as in the proof of Proposition~\ref{prop:lc-necessity}, now using the
height bound of Lemma~\ref{lem:height} for mixed boxes, and stop it after
$c\,\pi(I)^{c'}(I+q+1)^C$ bit operations, where $f(3,\kappa)\le
c\,\kappa^{c'}$ is the polynomial bound just noted. Under the promise the
algorithm has stopped by then with a certified interval
$[\beta,U]\ni\OPT$ of width at most $2^{-q}$, and the comparison of the
previous proof decides the question. If the clock expires, answer
arbitrarily; this is allowed outside the promise. The second statement
follows because a reduction whose image satisfies the promise would place
an NP-hard problem in $\mathrm P$.
\end{proof}

\begin{remark}[Exponential dependence on the bag size]\label{rem:lc-eth}
The factor exponential in $p$ cannot be avoided either, even where no
growth premise is needed. For a graph $G$ on $n$ vertices, the binary
quadratic program
$\min\{-\sum_ix_i+2\sum_{\{i,j\}\in E(G)}x_ix_j:\ x\in\{0,1\}^n\}$ has
optimal value equal to minus the independence number of $G$: if a selected vertex
has $k\ge1$ selected neighbours, deselecting it changes the objective by
$1-2k<0$. All diagonal entries are zero, so the problem lies in the
coordinatewise concave case, and a single bag containing all variables is
a valid decomposition with $p=n$. An algorithm for this class running in
time $2^{o(p)}\poly(I)$ would therefore find a maximum independent set in
time $2^{o(n)}$, which the exponential time hypothesis excludes
\cite{ImpagliazzoPaturiZane2001}. For this binary special case the
standard dynamic program takes time $2^{p}\poly(I)$, and Lokshtanov, Marx
and Saurabh show that maximum independent set admits no
$(2-\epsilon)^{\mathrm{tw}(G)}n^{O(1)}$ algorithm unless the strong
exponential time hypothesis fails \cite{LokshtanovMarxSaurabh2018}.
\end{remark}

\subsection{Related work}\label{sec:lc-related}

\paragraph{Sparse box QP, treewidth, and conditioning.}
For binary quadratic and pseudo-Boolean optimization, bounded treewidth
gives polynomial algorithms by nonserial dynamic programming
\cite{BerteleBrioschi1972,CramaHansenJaumard1990,Dechter1999} and exact
lift-and-project or LP formulations of size exponential only in the width
\cite{BienstockOzbay2004,WainwrightJordan2004,Laurent2009}; sparse
moment--SOS hierarchies exploit the same structure for continuous
polynomial problems \cite{WakiEtAl2006,Lasserre2006}. For continuous box QP
the situation is different. Del Pia and Khajavirad give an exact strongly
polynomial dynamic program on forests and prove strong NP-hardness at
treewidth two, even with bounded integral coefficients; they also show
that quartic minimization is strongly NP-hard on paths
\cite[Theorems~1--3]{DelPiaKhajavirad2026}. Their tractable classes of
unbounded treewidth, and those of Khajavirad
\cite{Khajavirad2026Poly}, separate variables with nonpositive diagonal
coefficients, which can be taken binary, from the continuous part. For
convex quadratic objectives with indicator variables, Bhathena et al.\
give an exact parametric dynamic program on trees
\cite{BhathenaEtAl2026Trees} and, on bounded-treewidth graphs, one whose
complexity depends on treewidth, volume growth, a margin and conditioning
\cite{BhathenaEtAl2026Graphs}; see also \cite{BienstockChen2024}. For integer QP, treewidth must be combined with
bounded domains or coefficients \cite{EibenEtAl2019,GanianOrdyniak2018,Lokshtanov2015}.
Approximation results with polynomial dependence on $1/\varepsilon$ are
known for sparse polynomial problems with constraints
\cite[Theorem~4]{BienstockMunoz2018} and for QP and MIQP with a fixed number of negative eigenvalues (or
fixed rank) and, for MIQP, a fixed number of integer variables
\cite{Vavasis1992,DelPia2023,DelPia2026Jacobi};
several of these works explain why the dependence on $1/\varepsilon$
cannot be made polylogarithmic in their settings unless
$\mathrm P=\mathrm{NP}$ \cite{BienstockMunoz2018,Vavasis1992,DelPia2026Jacobi}.
Proposition~\ref{prop:lc-necessity} is the analogous statement for
treewidth: some conditioning parameter is necessary. To our knowledge, no
earlier result gives a running time of the form
$f(p,\kappa)\poly(I+q)$, or an exact algorithm with time
$f(p,\kappa)\poly(I)$, for nonconvex or mixed-integer box QPs.

\paragraph{Grid and underestimator lower bounds.}
Lower bounds from function values on a grid are classical for Lipschitz
functions \cite{Piyavskii1972,Shubert1972,Nesterov2004}; with second-order
information the grid error is quadratic in the mesh
\cite{DeKlerkLasserreLaurentSun2017}. Separable quadratic perturbations
that convexify a function ($\alpha$BB) have maximum separation
$\frac14\sum_i\alpha_i(u_i-l_i)^2$
\cite{MaranasFloudas1994,AndroulakisMaranasFloudas1995,AdjimanEtAl1998};
perturbations in the opposite direction make a function edge-concave, so
that its convex envelope is vertex polyhedral
\cite{Tardella2004,MeyerFloudas2005,Hasan2018,MisenerFloudas2012}. Bajaj
and Hasan use an upper bound on the diagonal Hessian entries to obtain a
valid lower bound on a box from function values at its $2^n$ vertices
\cite{BajajHasan2020}. Piecewise-linear MIP relaxations of univariate
squares, such as sawtooth relaxations with error $2^{-2L-2}$ at depth $L$,
apply the same second-order error termwise inside a MIP
\cite{BeachHildebrandHuchette2022,BeachEtAl2024,DongLuo2018}.
Proposition~\ref{prop:lc-cellwise} shows that our corrected grid bound is
exactly the best vertex bound over all cells of a product grid, and that
the unary form of the correction lets one tree computation evaluate it
without enumerating cells.

\paragraph{Dynamic programming with coarse states for continuous variables.}
Coarse-to-fine dynamic programming replaces groups of states by
superstates with optimistic arc costs, so that each coarse problem is a
lower bound; it refines along the current optimal path and, for continuous
states, converges to the unique optimum \cite{Raphael2001}; see also
\cite{FelzenszwalbMcAllester2007}. Interval convex max-product bounds
potentials on interval cells to obtain a discrete lower bound for
continuous Markov random fields and refines the intervals
\cite{PengEtAl2011}. Interval dynamic programming for optimal power flow on
trees returns approximately feasible superoptimal solutions with a number
of bound-propagation calls polynomial in $1/\varepsilon$
\cite[Theorem~2]{DvijothamEtAl2017}. Continuous DCOP algorithms
\cite{HoangEtAl2020,TroullinosEtAl2022}, variable-resolution
discretization in control \cite{MunosMoore2002}, corridor methods
\cite{HeidariEtAl1971} and adaptive two-level dynamic programming
\cite{LeaverFayKuhlmanSnoeyink2005} refine discretizations adaptively.
None of these works gives a state count independent of the accuracy.

\paragraph{Filtering.}
Removing a region whose relaxation value exceeds an incumbent is
optimality-based bounds tightening, probing and branch-and-reduce in global
optimization \cite{RyooSahinidis1996,BelottiEtAl2009,GleixnerEtAl2017,PuranikSahinidis2017},
cost-based domain filtering in constraint programming
\cite{FocacciLodiMilano1999}, and dead-end elimination in protein design,
including provable pruning of continuous rotamer regions
\cite{DesmetEtAl1992,GainzaRobertsDonald2012}. Exact min-marginals
(max-marginals) of tree-structured problems are computed by two passes of
message passing \cite{Dawid1992,AjiMcEliece2000,WainwrightJaakkolaWillsky2005,KohliTorr2006},
and thresholding them prunes states in structured prediction cascades
\cite{WeissTaskar2010}. Our filter is an instance of this principle; its
role here is that, under quadratic growth, the retained coordinate hulls
have radius proportional to the current mesh. A proof of optimality that
records every pruning decision is also standard in branch and bound
\cite{CheungGleixnerSteffy2017}.

\paragraph{Accuracy-independent work under growth.}
In branch and bound, second-order convergent lower bounds and a
nondegenerate minimizer make the number of boxes near the minimizer
independent of the tolerance (the cluster problem); this count can still be
exponential in the dimension and in the conditioning of the Hessian
\cite{DuKearfott1994,Neumaier2004,WechsungSchaberBarton2014,KannanBarton2017}.
Similarly, optimistic optimization attains an exponential rate when the
growth order at the optimum matches the assumed smoothness, without knowing
the smoothness \cite{Munos2011}, and for Lipschitz functions whose
$\varepsilon$-optimal sets have measure of order $\varepsilon^n$ a
branch-and-bound search needs only $O(\log(1/\varepsilon))$ function
values \cite{Perevozchikov1990}. Proximity theorems and scaling give algorithms whose dependence on the
domain size is logarithmic for linear and separable convex integer programs
\cite{CookGerardsSchrijverTardos1986,HochbaumShanthikumar1990,GranotSkorinKapov1990,HunkenschroderEtAl2026},
and sparse integer programs, also with separable convex objectives, are
fixed-parameter tractable in a sparsity measure and the coefficient size
\cite{EisenbrandEtAl2025}.
Quadratic growth and error bounds give linear convergence of local methods
\cite{LuoTseng1993,KarimiNutiniSchmidt2016,DrusvyatskiyLewis2018,NecoaraNesterovGlineur2019},
and restarts remove the need to know the growth constant
\cite{RouletDAspremont2020}. Our contraction argument is a global,
non-asymptotic version of the cluster phenomenon in which filtering is
coordinatewise; this replaces the exponential dependence on the dimension
by $\log n$ per coordinate and $\max_i|G_i|^p$ per bag.

\paragraph{Exact rational output.}
A rational QP that attains its minimum has a global optimizer of polynomial
encoding length \cite{Vavasis1990,DelPiaDeyMolinaro2017}; the proof of
Lemma~\ref{lem:height} follows Vavasis's argument. Rounding an approximate
solution to an exact rational one by continued fractions is classical
\cite[Section~5.1]{GrotschelLovaszSchrijver1988}, \cite{KozlovTarasovKhachiyan1980}.
Uniqueness of a QP minimizer implies quadratic growth by second-order
optimality theory for quadratic programs, in which second-order
necessary and sufficient conditions coincide \cite{Contesse1980}, together
with compactness; see also \cite{BonnansIoffe1995}; set-valued error bounds
for quadratic systems are due to Luo and Sturm \cite{LuoSturm2000}. For
polynomials of degree three or more, exact output is limited by irrational
optimizers and encoding-size phenomena \cite{BienstockDelPiaHildebrand2023}.
